\section{Introduction}
A representation changes which distinctions an agent can use. A repair can
improve task performance, alter the label of the representation currently
installed, and change the observations available for evaluating earlier
decisions. These effects create two questions that should be studied together
without identifying their objectives: how much evidence is needed to certify
an old representation when repair is allowed, and when should an agent
intervene under a finite task deadline?

For old-label certification, repairing first can change which pre-repair
alternatives must be ruled out. Suppose three environments have old labels
$(0,0,1)$, and a repair reveals whether the environment is in $\{1\}$ or
$\{2,3\}$. In environment 1, it may be efficient to rule out environment 2
and then use the decoder $(0,1)$. The unresolved environment 3 has the
opposite old label, but the post-repair observation can distinguish it from
environment 1. Certifying the final old label before any repair would instead
need to rule out environment 3 directly.

For current-label task control, the relevant decision can differ even without
a high-confidence constraint. A Bayes declaration compares the two costs of
misclassifying the installed representation. An intervention compares task
losses, its own price, remaining time, and future information. These comparisons
can favor different actions at the same belief. A noisy verification option can
also split the optimal repair region into disconnected intervals.

We formalize both questions in finite known models with one joint
observation--state convention and explicit loss accounting. All declared
theorems include proofs. The principal results are:
\begin{enumerate}
\item A reduction from an ordered repair experiment to a finite
multiple-correct-answer problem, with optimal old-label certification cost
$g_i\log(1/\delta)/D_i$, a common attaining policy, and fully charged deadline
tails (Theorem~\ref{thm:main}).
\item A factorization of $D_i$ across post-repair classes and a polynomial
eligibility test for the same likelihood-ratio policy
(Theorems~\ref{thm:factorization} and \ref{thm:polynomial}).
\item A finite-deadline information bound and a matching shrinking-channel
regime in which the joint coefficient returns to the certification-before-repair
coefficient, despite distinct post laws at every confidence level
(Theorems~\ref{thm:deadline-lower} and \ref{thm:collapse}).
\item For finite-horizon Bayes control, an exact declaration-gate cost identity,
a bound on the number of repair-region components, and a full-support TP2
example attaining two components (Theorems~\ref{thm:gate},
\ref{thm:geometry}, and \ref{thm:counter}).
\item Exact no-information gate costs and a constructive logarithmic gate-cost
upper bound with repeatable identifiable verification
(Proposition~\ref{prop:noinfo} and Theorem~\ref{thm:log}), with reproducible
finite computations under each stated contract.
\end{enumerate}

The common message is precise: evidence required for a final assertion and
evidence sufficient to select a useful intervention need not coincide.
The two policy restrictions in this paper are different. Certification before
repair commits to a fixed old-label answer under a uniform error constraint.
A declaration gate restricts interventions according to a Bayes decision
boundary about the current representation. We compare policies only within
one objective and one set of physical kernels at a time.

\paragraph{Scope and contribution.}
The statistical max--min answer complexity, belief-state recursion, and
information-comparison tools are established ideas. Our focus is their explicit
repair interpretation, the classwise factorization, finite-deadline boundary,
and repair-region structure in the stated models. No result assumes that
calling a state a representation automatically makes a new control theory.
Proposition~\ref{prop:realize} shows why arbitrary finite risk tables can fit
the adequacy definition. Open-ended representation discovery, unknown model
learning, and arbitrary repeated repair remain outside the proved scope.

\section{Related work}
\label{sec:related}
The closest statistical comparison is pure exploration with multiple correct
answers \citep{degenne2019}. Its answer-dependent alternatives become the
invalidating environments of a post-repair decoder. With only one pre channel,
there is no experiment-allocation optimization. The coefficient's max--min
form is therefore not claimed as a new principle. The ordered-process
reduction and deadline accounting are established explicitly, while the
disjoint class coordinates yield the polynomial implementation in
Section~\ref{sec:decoder-structure}.

Sequential change of measure connects sample counts and information
\citep{kaufmann2016}; non-uniform costs and controlled sensing have also been
studied \citep{nitinawarat2015}. We use finite likelihood-ratio martingales and
bounded-increment concentration \citep{hoeffding1963}. The finite-deadline
bound accounts for information from both phases and exposes a nonuniform
limit of the free-post-information coefficient.

Finite-horizon POMDP belief states and piecewise-linear value functions are
classical \citep{smallwood1973}. Their concavity supplies the repair-region
component argument; the Bellman equation itself is not a contribution.
Inspection and maintenance already combine observation and intervention
\citep{morato2022,zhang2023}. Structural POMDP results require conditions beyond
an informative ordered channel \citep{krishnamurthy2018}. Our TP2 example shows
that this channel condition alone is insufficient for a connected repair
region, and does not contradict theorems with additional cost and transition
conditions. Pure-information monotonicity follows from Blackwell comparison
\citep{blackwell1953}.

Model editing provides a broader example of changing stored computational
content \citep{meng2022memit}, while learning to defer studies a predictor's
use of an expert \citep{mozannar2020defer}. These motivate repair and fallback,
but we neither model a specific editing method nor claim empirical gains for
a deployed language-model system. The fixed-representation M1 manuscript
\citep{huang2026} motivates the adequacy question. The present proofs are
self-contained; its fixed-confidence objective is not silently identified
with a finite Bayes horizon.

\paragraph{Notation.}
All logarithms are natural. For finite laws,
$\KL(P\Vert Q)=\sum_yP(y)\log(P(y)/Q(y))$.
We use $\min\varnothing=+\infty$ and $c/(+\infty)=0$.
A uniform error guarantee belongs to one implementable policy used in every
environment. An environment-indexed infimum still ranges over this policy
class; the controller is never given the true index.
